\section{Discussion}

The value of this extension is not that it directly prescribes field policy. The value is that it offers a controlled way to compare institutional arrangements that correspond to recognizable governance choices in renewable commons. In that sense, the paper is intended as a decision-support and theory-building tool rather than a literal policy simulator.

This positioning matters for both Cooperative AI and commons scholarship. For Cooperative AI, the paper studies how collective systems remain stable under strategic pressure when governance is imperfect. For commons research, it provides a controlled way to compare central, local, and hybrid institutional arrangements under forms of strategic noncompliance that are difficult to isolate in the field.

Future work can strengthen the outward bridge through closer calibration, additional empirical domain anchors, or collaborations with researchers working directly on fisheries, irrigation, or forest governance. The present contribution is a literature-backed institutional benchmark that remains grounded in the same research line while improving interpretability and relevance.
